
\setcounter{theorem}{\numexpr\getrefnumber{th:expoHD-cliff}-1\relax}


\begin{theorem}
	The {\oneoneOPTIA} using \IPHfcm{} with \linHD\ (factor $c<\tfrac{1+4\varepsilon}{3-4\varepsilon}$) and ageing parameter $\tau=\Omega(n^{1+\epsilon'})$
	optimises \cliff{} with $d<n(\tfrac{1}{4}-\epsilon)$ in expected
	\[
	O\!\Big(\frac{n^{5}\log(n/d)}{d^{3}} \;+\; n^{2}\log d\Big)
	\]
	fitness function evaluations.
\end{theorem}

\textcolor{blue}{In order to prove the above theorem we will first establish the supporting lemmata.}


\setcounter{lemma}{\getrefnumber{lem:can-start-omega-d3}-1\relax}


\begin{lemma}[Canonical start from an ageing cliff-top]
	Assume the process is at the cliff-top (zero-count $d$) when ageing is triggered, under \IPHfcm{} with mutation potential $M(x)=\max\{1,\lfloor c\,\HD(x,x^\top)\rfloor\}$ for some $c\in(0,1)$, the reference-update rule (“update when a strictly better solution is sampled or a best-fitness solution dies”), and the convention that \emph{ageing resets only on strict improvement}. Then, within at most three iterations, the process reaches the canonical second-slope start
	\[
	r_0=2\qquad\text{and}\qquad \HD(r_0)=2
	\]
	with probability at least $\displaystyle \frac{3}{256}\,\bigl(\tfrac{d}{n}\bigr)^{3}=\Omega\!\bigl((d/n)^3\bigr)$.
\end{lemma}

%\begin{proof}
%	We exhibit a three-step event whose occurrence yields the stated start and whose probability is $\Omega((d/n)^3)$.
%	
%	\emph{(A) Equal-fitness lock at the ageing boundary.} In the ageing iteration, require the first two flips to be $1\!\to 0$ then $0\!\to 1$ (without replacement); this has probability $\frac{n-d}{n}\cdot\frac{d}{n-1}$ and produces an \emph{equal-fitness} offspring accepted by FCM (first return to $\ge$ parent). The offspring inherits the parent’s age (no improvement), so both parent and offspring face the ageing coin and die independently with probability $1/2$. Requiring “parent dies, offspring survives’’ contributes a factor $1/4$, and by the update rule the reference locks to the last cliff-top (the parent that died). The surviving equal-fitness offspring differs from this locked reference in exactly two bits, hence $\HD=2$ and, since $c<1$, the next iteration has $M=1$.
%	
%	\emph{(B) Create the cliff-bottom by $M{=}1$ exhaustion and kill the current cliff-top.} With $M=1$, exactly one flip is performed and FCM does not accept worse-than-parent, so the algorithm returns the single-flip string after potential exhaustion. Requiring a $0\!\to 1$ flip occurs with probability at least $d/n$ and returns the \emph{cliff-bottom} ($n-d+1$ ones), at Hamming distance $1$ from the locked reference. Ageing does not reset (no improvement), so both current parent (equal-fitness cliff-top) and the returned offspring face the coin; requiring “parent dies, offspring survives’’ contributes another factor $1/4$ and preserves the locked reference at the last cliff-top.
%	
%	\emph{(C) One accepted improvement on the second slope.} In the following iteration, require the first flip to be $0\!\to 1$; this occurs with probability at least $(d-1)/n$ and is \emph{accepted} by FCM (strict improvement from the cliff-bottom), producing a point with $n-d+2$ ones. Relative to the locked cliff-top reference, the Hamming distance is now $2$, and exactly two successful phases have completed: $r_0=2$, $\HD(r_0)=2$.
%	
%	Multiplying the lower bounds yields
%	\[
%	\Big(\tfrac{n-d}{n}\cdot\tfrac{d}{n-1}\cdot\tfrac14\Big)\cdot
%	\Big(\tfrac{d}{n}\cdot\tfrac14\Big)\cdot
%	\Big(\tfrac{d-1}{n}\Big).
%	\]
%	For $d\le(\tfrac14-\varepsilon)n$ we have $(n-d)/n\ge 3/4$, and for $d\ge 2$ we have $\tfrac{d}{n-1}\ge\tfrac{d}{n}$ and $\tfrac{d-1}{n}\ge \tfrac12\cdot\tfrac{d}{n}$. Hence the product is at least
%	\[
%	\frac{3}{256}\,\Big(\frac{d}{n}\Big)^{\!3},
%	\]
%	which proves the claim.
%\end{proof}
%
%\setcounter{lemma}{\getrefnumber{lem:xi-tail}-1\relax}
%
%\begin{lemma}[Exponential tail for catch-up length]
%	Fix a state on the second slope with $z$ zero-bits (so $z\le d<(\tfrac14-\varepsilon)n$) and let $p:=z/n$.
%	During one hypermutation, let $\xi_z$ be the number of $1\!\to 0$ flips that occur \emph{before} the process first returns to fitness $\ge$ parent (i.e., before the first ``catch-up'' to deficit $0$). Then for every $k\ge 1$,
%	\[
%	\Pr(\xi_z\ge k)\ \le\ \Big(\tfrac{2z}{n}\Big)^k.
%	\]
%	In particular, the event that the Hamming distance increases by at least $2k$ before stopping (equivalently, the first catch-up occurs at time $\ge 2k$) has probability at most $(2z/n)^k$.
%\end{lemma}
%
%\begin{proof}
%	If the first flip is a $0\!\to 1$ move, the offspring is strictly better and the hypermutation stops immediately, yielding $\xi_z=0$. Hence the event $\{\xi_z\ge 1\}$ implies the first flip was a $1\!\to 0$ move. From that point on the process carries a \emph{deficit} $D_t:=\#(1\!\to 0)-\#(0\!\to 1)$ that starts at $D_1=1$ and decreases by $1$ on each $0\!\to 1$ flip and increases by $1$ on each $1\!\to 0$ flip. The hypermutation stops at the \emph{first} time $t$ with $D_t=0$ (a catch-up). Thus the event $\{\xi_z\ge k\}$ is exactly the event that the first return to $0$ happens no earlier than time $2k$.
%	
%	We upper bound this probability by a simple ``coarse blocking’’ argument that uses only that $z$ is small and flips are without replacement.
%	
%	\smallskip
%	\emph{Step 1 (Each required zero is never more likely than $2p$).}
%	Condition on any history of length $t\le 2k-1$ that keeps $D_t\ge 1$ (so no catch-up yet) and in which at most $t$ distinct positions have been flipped. Among the $n-t$ remaining positions there are at most $z$ zeros. Hence the conditional probability that the \emph{next} flip is a $0\!\to 1$ move is at most
%	\[
%	\frac{z}{\,n-t\,}\ \le\ \frac{z}{\,n-(2k-1)\,}\ \le\ \frac{z}{(1/2+2\varepsilon)n}\ \le\ 2p,
%	\]
%	because $2k\le 2d<(\tfrac12-2\varepsilon)n$. In words: throughout any run that lasts up to $2k$ flips, a zero-flip is never more likely than $2p$.
%	
%	\smallskip
%	\emph{Step 2 (Decompose a catch-up of length $\ge 2k$ into $k$ gated zero-flips).}
%	If the first catch-up occurs at some time $\ge 2k$, then along the way the trajectory must realize \emph{at least $k$} distinct $0\!\to 1$ flips while $D>0$ (each such flip reduces the deficit by $1$ but, by assumption, does not yet hit $0$ before the $k$-th such zero). Expose these $k$ zero-flips in order and ignore everything else (the interspersed $1\!\to 0$ flips only make catch-up \emph{harder} and so can be dropped for an upper bound). By Step~1, each of these $k$ zeros occurs with conditional probability at most $2p$, uniformly over the entire run while $D>0$. Therefore, by the chain rule,
%	\[
%	\Pr(\text{first catch-up occurs at time }\ge 2k)\ \le\ (2p)^k.
%	\]
%	
%	
%	Since the left-hand side is exactly $\Pr(\xi_z\ge k)$, the claim follows.
%\end{proof}

\begin{proof}
	We exhibit a consecutive three-step event which happens with probability $\Omega((d/n)^3)$ and yields the stated starting 
	position on the second slope.
	
	\emph{(A) $M=1$ at the ageing boundary.} In the ageing iteration we require that the first two flips are $1\!\to 0$ then $0\!\to 1$ without replacement. This has probability $\frac{n-d}{n}\cdot\frac{d}{n-1}$ and produces an \emph{equal-fitness} offspring accepted by FCM (first return to $\ge$ parent). The offspring inherits the parent’s age (no improvement), so both parent and offspring face the ageing coin and die independently with probability $1/2$. Requiring “parent dies, offspring survives’’ contributes a factor $1/4$, and by the update rule the reference locks to the last cliff-top (the parent that died). The surviving equal-fitness offspring differs from this locked reference in exactly two bits, hence $\HD=2$ and, since $c<1$, the next iteration has $M=1$.
		
	\emph{(B) Create the cliff-bottom individual and kill the current cliff-top.} With $M=1$, exactly one flip is performed and FCM does not accept worse-than-parent, so the algorithm returns the single-flip string after potential exhaustion. A $0\!\to 1$ flip occurs with probability at least $d/n$ and returns the \emph{cliff-bottom} ($n-d+1$ ones), at Hamming distance $1$ from the locked reference. Ageing does not reset (no improvement), so both current parent (equal-fitness cliff-top) and the returned offspring face the coin; requiring “parent dies, offspring survives’’ contributes another factor $1/4$ and preserves the locked reference at the last cliff-top.
	
	\emph{(C) One accepted improvement on the second slope.} In the following iteration with probability at least $(d-1)/n$ the first bit flip iss a $0\!\to 1$ flip, and the offspring is \emph{accepted} by FCM (strict improvement from the cliff-bottom), producing a point with $n-d+2$ ones. Relative to the locked cliff-top reference, the Hamming distance is now $2$, and exactly two successful phases have completed: $r_0=2$, $\HD(r_0)=2$.
	
	Multiplying the lower bounds yields
	\[
	\Big(\tfrac{n-d}{n}\cdot\tfrac{d}{n-1}\cdot\tfrac14\Big)\cdot
	\Big(\tfrac{d}{n}\cdot\tfrac14\Big)\cdot
	\Big(\tfrac{d-1}{n}\Big).
	\]
	For $d\le(\tfrac14-\varepsilon)n$ we have $(n-d)/n\ge 3/4$, and for $d\ge 2$ we have $\tfrac{d}{n-1}\ge\tfrac{d}{n}$ and $\tfrac{d-1}{n}\ge \tfrac12\cdot\tfrac{d}{n}$. Hence the product is at least
	\[
	\frac{3}{256}\,\Big(\frac{d}{n}\Big)^{\!3},
	\]
	which proves the claim.
\end{proof}

\setcounter{lemma}{\getrefnumber{lem:xi-tail}-1\relax}

\begin{lemma}[Exponential tail for catch-up length]
	Fix a state on the second slope with $z$ zero-bits (so $z\le d<(\tfrac14-\varepsilon)n$) and let $p:=z/n$.
	During one hypermutation, let $\xi_z$ be the number of $1\!\to 0$ flips that occur \emph{before} the process first returns to fitness $\ge$ parent (i.e., before the first ``catch-up'' to deficit $0$). Then for every $k\ge 1$,
	\[
	\Pr(\xi_z\ge k)\ \le\ \Big(\tfrac{2z}{n}\Big)^k.
	\]
	In particular, the event that the Hamming distance increases by at least $2k$ before stopping (equivalently, the first catch-up occurs at time $\ge 2k$) has probability at most $(2z/n)^k$.
\end{lemma}

%\begin{proof}
%	If the first flip is a $0\!\to 1$ move, the offspring is strictly better and the hypermutation stops immediately, yielding $\xi_z=0$. Hence the event $\{\xi_z\ge 1\}$ implies the first flip was a $1\!\to 0$ move. From that point on the process carries a \emph{deficit} $D_t:=\#(1\!\to 0)-\#(0\!\to 1)$ that starts at $D_1=1$ and decreases by $1$ on each $0\!\to 1$ flip and increases by $1$ on each $1\!\to 0$ flip. The hypermutation stops at the \emph{first} time $t$ with $D_t=0$ (a catch-up). Thus the event $\{\xi_z\ge k\}$ is exactly the event that the first return to $0$ happens no earlier than time $2k$.
%	
%	We upper bound this probability by a simple ``coarse blocking’’ argument that uses only that $z$ is small and flips are without replacement.
%	
%	\smallskip
%	\emph{Step 1 (Each required zero is never more likely than $2p$).}
%	Condition on any history of length $t\le 2k-1$ that keeps $D_t\ge 1$ (so no catch-up yet) and in which at most $t$ distinct positions have been flipped. Among the $n-t$ remaining positions there are at most $z$ zeros. Hence the conditional probability that the \emph{next} flip is a $0\!\to 1$ move is at most
%	\[
%	\frac{z}{\,n-t\,}\ \le\ \frac{z}{\,n-(2k-1)\,}\ \le\ \frac{z}{(1/2+2\varepsilon)n}\ \le\ 2p,
%	\]
%	because $2k\le 2d<(\tfrac12-2\varepsilon)n$. In words: throughout any run that lasts up to $2k$ flips, a zero-flip is never more likely than $2p$.
%	
%	\smallskip
%	\emph{Step 2 (Decompose a catch-up of length $\ge 2k$ into $k$ gated zero-flips).}
%	If the first catch-up occurs at some time $\ge 2k$, then along the way the trajectory must realize \emph{at least $k$} distinct $0\!\to 1$ flips while $D>0$ (each such flip reduces the deficit by $1$ but, by assumption, does not yet hit $0$ before the $k$-th such zero). Expose these $k$ zero-flips in order and ignore everything else (the interspersed $1\!\to 0$ flips only make catch-up \emph{harder} and so can be dropped for an upper bound). By Step~1, each of these $k$ zeros occurs with conditional probability at most $2p$, uniformly over the entire run while $D>0$. Therefore, by the chain rule,
%	\[
%	\Pr(\text{first catch-up occurs at time }\ge 2k)\ \le\ (2p)^k.
%	\]
%	
%	
%	Since the left-hand side is exactly $\Pr(\xi_z\ge k)$, the claim follows.
%\end{proof}
%

\begin{proof}
	If the first flip is a $0\!\to 1$ move, the offspring is strictly better and the hypermutation stops immediately, yielding $\xi_z=0$. Hence the event $\{\xi_z\ge 1\}$ implies that the first flip was a $1\!\to 0$ move. From that point on the process carries a \emph{deficit} of one-bits $D_t:=\#(1\!\to 0)-\#(0\!\to 1)$ that starts at $D_1=1$ and decreases by $1$ on each $0\!\to 1$ flip and increases by $1$ on each $1\!\to 0$ flip. The hypermutation stops at the \emph{first} time $t$ with $D_t=0$ (a catch-up). Thus the event $\{\xi_z\ge k\}$ is exactly the event that the first return to $0$ happens no earlier than time $2k$.
	
	We bound this probability from above by a simple ``coarse blocking’’ argument that uses only that $z$ is small and that flips occur are without replacement.
	
	\smallskip
	\emph{Step 1 (Each required zero-flip is never more likely than $2p$).}
	Condition on any history of length $t\le 2k-1$ that keeps $D_t\ge 1$ (so no catch-up yet) and in which at most $t$ distinct positions have been flipped. Among the $n-t$ remaining positions there are at most $z$ zeros. Hence the conditional probability that the \emph{next} flip is a $0\!\to 1$ move is at most
	\[
	\frac{z}{\,n-t\,}\ \le\ \frac{z}{\,n-(2k-1)\,}\ \le\ \frac{z}{(1/2+2\varepsilon)n}\ \le\ 2p,
	\]
	because $2k\le 2d<(\tfrac12-2\varepsilon)n$. In other words, throughout any hypermutation that lasts up to $2k$ flips, a zero-flip is never more likely than $2p$.
	
	\smallskip
	\emph{Step 2 (Decompose a catch-up of length $\ge 2k$ into $k$ gated \textcolor{purple}{?} zero-flips).}
	If the first catch-up occurs at some time $\ge 2k$, then along the way the trajectory must realize \emph{at least $k$} distinct $0\!\to 1$ flips while $D>0$ (each such flip reduces the deficit by $1$ but, by assumption, does not yet hit $0$ before the $k$-th such zero\new{-flip}). Expose these $k$ zero-flips in order and ignore everything else (the interspersed $1\!\to 0$ flips only make catch-up \emph{harder} and so can be dropped for an upper bound). By Step~1, each of these $k$ zero\new{-flips} occurs with conditional probability at most $2p$, uniformly over the entire run while $D>0$. Therefore, by the chain rule,
	\[
	\Pr(\text{first catch-up occurs at time }\ge 2k)\ \le\ (2p)^k.
	\]
	
	
	Since the left-hand side is exactly $\Pr(\xi_z\ge k)$, the claim follows.
\end{proof}




%\begin{lemma}[Per-attempt MGF]\label{lem:attempt-mgf}
%	Fix a state on the second slope with $z$ zero-bits and write $p:=z/n$.
%	Within a single hypermutation attempt at zero-count $z$, let $\Delta_z$ denote the Hamming-distance gain accumulated \emph{until the attempt stops} (by $\mathrm{FCM}(\ge)$ or by potential exhaustion).
%	Then for every $0<\lambda<\tfrac12\ln\!\frac{1}{2p}$,
%	\[
%	\EE\!\big[e^{\lambda(\Delta_z-1)}\big]\ \le\ \frac{1}{\,1-2p\,e^{2\lambda}\,}.
%	\]
%	In particular, since $z\le d$ and $\rho:=\tfrac{2d}{n}$, for $0<\lambda<\tfrac12\ln\!\tfrac{1}{\rho}$,
%	\[
%	\EE\!\big[e^{\lambda(\Delta_z-1)}\big]\ \le\ \frac{1}{\,1-\rho\,e^{2\lambda}\,}.
%	\]
%\end{lemma}
%
%\begin{proof}
%	Let $\xi_z$ be the number of $1\!\to 0$ flips that occur before the first catch-up to fitness $\ge$ parent (as in Lemma~\ref{lem:xi-tail}).
%	Within any attempt we have the crude envelope
%	\[
%	\Delta_z\ \le\ 1+2\,\xi_z,
%	\]
%	because each additional $1\!\to 0$ incurred while below the parent must eventually be neutralised by a $0\!\to 1$ to reach equality (or else the attempt stops earlier, which only decreases $\Delta_z$).
%	Hence
%	\[
%	e^{\lambda(\Delta_z-1)}\ \le\ e^{2\lambda\,\xi_z}.
%	\]
%	For $s:=e^{2\lambda}>1$ we therefore get
%	\[
%	\EE\!\big[e^{\lambda(\Delta_z-1)}\big]\ \le\ \EE[s^{\xi_z}].
%	\]
%	By Lemma~\ref{lem:xi-tail}, $\Pr(\xi_z\ge k)\le (2p)^k$ for all $k\ge 1$.
%	Using the tail-sum inequality for nonnegative integer random variables,
%	\begin{align*}
	%	\EE[s^{\xi_z}]
	%	\ &=\ 1+\sum_{k\ge 1}\big(s^{k}-s^{k-1}\big)\Pr(\xi_z\ge k)
	%	\ \\&\le\ 1+\sum_{k\ge 1}s^{k}\Pr(\xi_z\ge k)
	%	\ \le\ 1+\sum_{k\ge 1}\big(2p\,s\big)^{k}
	%	\ \\ &=\ \frac{1}{\,1-2p\,s\,},
	%	\end{align*}
%	provided $2ps<1$, i.e., $0<\lambda<\tfrac12\ln\!\tfrac{1}{2p}$. This proves the first claim.
%	If $z\le d$ then $2p\le\rho$, and since the RHS increases with $p$, we obtain
%	\(
%	\EE[e^{\lambda(\Delta_z-1)}]\le (1-\rho e^{2\lambda})^{-1}
%	\)
%	for any $0<\lambda<\tfrac12\ln(1/\rho)$.
%\end{proof}

\begin{lemma}[Phase MGF, refined]\label{lem:phase-mgf}
	Let a phase be the time spent at zero-count $z$ on the second slope until the next strict improvement to $z-1$.
	Write $p:=z/n$ and $\rho:=2d/n$ (so $2p\le \rho \le \tfrac12-2\varepsilon$).
	Let $S_z$ be the total HD added by all \emph{failed} attempts in the phase.
	Since the improving attempt must flip a zero at the first flip  we have $\Delta^{\mathrm{succ}}_z=1$. Hence
	\[
	Y_z\ :=\ \textcolor{purple} {S_z +(\Delta^{\mathrm{succ}}_z-1)}\ =\ S_z\ \ge 0.
	\]
	Then, for every $0<\lambda<\tfrac12\ln\frac{1}{2\rho}$,
	\[
	\EE\!\big[e^{\lambda Y_z}\big]\ \le\ K(\lambda)\ :=\ \frac{1-\rho}{\,1-\rho\,e^{2\lambda}\,}.
	\]
\end{lemma}

\begin{proof}
	For a single \emph{failed} attempt starting the hypermutation with $z$ zeroes, let $\xi$ count the number of $1\!\to 0$ flips while below the parent's fitness before the first catch-up.
	By Lemma~\ref{lem:xi-tail} and the change in Hamming distance $\Delta_{\mathrm{fail}}\le 2\xi$ with $\xi\ge 1$,
	\begin{align*}
		\EE\!\left[e^{\lambda \Delta_{\mathrm{fail}}}\right]
		&\le \sum_{k\ge 1}^{\infty} e^{2\lambda k}\,\Pr(\xi=k)
		\\ &\le\ \sum_{k\ge 1}^{\infty} \big(2p\,e^{2\lambda}\big)^{k}
		\\ &\textcolor{purple}{=}\ \frac{2p\,e^{2\lambda}}{\,1-2p\,e^{2\lambda}\,}:=:\ t.
	\end{align*}
	Let $F$ be the number of failed attempts before the first success in the phase. %with the convention
	%$\Pr(F=f)=(1-p_z)^f p_z$ for $f\in\{0,1,\ldots\}$ and $p_z\ge p=z/n$.
	Since $Y_z$ is the sum over failed attempts only (the success contributes $0$), conditioning stepwise on the phase history gives
	\begin{align*}
		\EE\!\big[e^{\lambda Y_z}\big]
		&=\EE\!\Big[\ \prod_{i=1}^{F} e^{\lambda \Delta_{\mathrm{fail}}^{(i)}}\ \Big]
		\ \le\ \EE\!\big[t^{F}\big].
	\end{align*}
	Using the geometric MGF for $t\in[0,1)$,
	\begin{align*}
		\EE\!\big[t^{F}\big]
		&=
		 \sum_{f\ge 0} (1-p)^f p\, t^f
		\ =\ \frac{p}{\,1-(1-p)t\,}.
	\end{align*}
	Substituting $t=\dfrac{2p\,e^{2\lambda}}{1-2p\,e^{2\lambda}}$ and simplifying,
	\begin{align*}
		\EE\!\big[e^{\lambda Y_z}\big]
		&\le \frac{p}{\,1-(1-p)\,\dfrac{2p\,e^{2\lambda}}{1-2p\,e^{2\lambda}}\,}
		\\ &=\ \frac{p\,(1-2p\,e^{2\lambda})}{\,1-2p\,e^{2\lambda}(2-p)\,}
		\ :=:\ K_z(\lambda).
	\end{align*}
	Since $p\le \rho/2$ and the map $p\mapsto K_z(\lambda)$ is increasing on $(0,\tfrac12)$ for fixed $\lambda\in(0,\tfrac12\ln(1/\rho))$, we obtain the uniform bound
	\begin{align*}
		\EE\!\big[e^{\lambda Y_z}\big]
		\ \le\ K_z(\lambda)\Big|_{p=\rho/2}=\frac{\tfrac{\rho}{2}(1-\rho\,e^{2\lambda})}{1-\rho e^{2\lambda}\,(2-\tfrac{\rho}{2})}
	\end{align*}
	Now, our claimed upper bound reduces to showing that a concave quadratic function of $u:=\rho e^{2\lambda} \in \left[0,\frac{1}{2}\right]$ is non-negative.
	\begin{align*}
		&\frac{\tfrac{\rho}{2}(1-\rho\,e^{2\lambda})}{1-\rho e^{2\lambda}\,(2-\tfrac{\rho}{2})}
		\  \le\ \frac{1-\rho}{\,1-\rho\,e^{2\lambda}\,}
		\ =:\ K(\lambda),
		\\&\Longleftrightarrow\quad
		\frac{\rho}{2}(1-u)^2 \le (1-\rho)\Bigl(1-\bigl(2-\tfrac{\rho}{2}\bigr)u\Bigr),\ \ u:=\rho\,e^{2\lambda},
		\\
		D(u)
		&:=(1-\rho)\Bigl(1-\bigl(2-\tfrac{\rho}{2}\bigr)u\Bigr)-\frac{\rho}{2}(1-u)^2
		\\&= \Bigl(1-\tfrac{3}{2}\rho\Bigr)+\Bigl(-2+\tfrac{7}{2}\rho-\tfrac{1}{2}\rho^2\Bigr)u-\tfrac{\rho}{2}u^2,
		\\
		D(0)&=1-\tfrac{3}{2}\rho\ \ge\ \tfrac14,\qquad
		D\!\Big(\tfrac12\Big)=\tfrac{\rho}{8}\bigl(1-2\rho\bigr)\ \ge\ 0
		\quad\Rightarrow\quad \\
		D(u)&\ge 0\ \text{ for }u\in[0,\tfrac12],
	\end{align*}
	
	
\end{proof}



\begin{lemma}[Buffer process and safety invariant]\label{lem:buffer}
	Let successful phases be indexed by $r=r_0,r_0+1,\dots$ and write
	$s_0:=\HD(r_0)-r_0$, $\theta:=\tfrac{1}{c}-1$, 
	$E_r:=\sum_{j=r_0+1}^{r} Y_{z_j}$, and 
	$B_r:=\theta r - s_0 - E_r$.
	Then for all $r\ge r_0:$  \\
	(i) $c\cdot\HD(r)\le r \iff B_r\ge 0$; \\
	(ii) $B_{r+1}-B_r=\theta-Y_{z_{r+1}}$; \\
	(iii) $B_{r_0}=\theta r_0 - s_0$. 
\end{lemma}

\begin{proof}
	For any $r\ge r_0$, the Hamming distance admits the exact decomposition
	\begin{align*}
		\HD(r)\ &=\ \underbrace{\HD(r_0)}_{=\,r_0+s_0}\ +\ \underbrace{(r-r_0)}_{\text{one per success}}\ +\ \underbrace{E_r}_{\text{ extra HD in failed attempts}}
		\ \\&=\ r+s_0+E_r,		
	\end{align*}
	
	since each strict improvement increases the HD by exactly $1$ and $E_r=\sum_{j=r_0+1}^r Y_{z_j}$ collects the extra HD from failed attempts (the successful attempt contributes $0$ because $\Delta_z^{\mathrm{succ}}=1$).
	
	\smallskip
	\emph{(i) Safety invariant.}
	Using $\theta=\tfrac{1}{c}-1$ and $B_r:=\theta r - s_0 - E_r$,
	\begin{align*}
		B_r\ge 0
		\ &\Longleftrightarrow\
		\theta r \ge s_0+E_r
		\ \Longleftrightarrow\
		c\theta r \ge c(s_0+E_r)
		\\ &\Longleftrightarrow\
		(1-c)r \ge c(s_0+E_r).	
	\end{align*}
	
	
	But $\HD(r)=r+s_0+E_r$, so the last inequality is equivalent to
	\[
	r \ \ge\ c\,(r+s_0+E_r)\ =\ c\,\HD(r),
	\]
	i.e., $c\,\HD(r)\le r$. Thus $c\,\HD(r)\le r \iff B_r\ge 0$.
	
	\smallskip
	\emph{(ii) One-phase update.}
	By definition,
	\begin{align*}
		B_{r+1}-B_r
		&= \bigl(\theta(r+1)-s_0-E_{r+1}\bigr) - \bigl(\theta r - s_0 - E_r\bigr)
		\\&= \theta - (E_{r+1}-E_r)
		= \theta - Y_{z_{r+1}}.
	\end{align*}
	
	
	
	\smallskip
	(iii) Trivial from the definition at $r=r_0$.
\end{proof}

% -------------------------------
% Step 1: Exponential drift envelope
% -------------------------------
\begin{lemma}[Global exponential drift envelope]\label{lem:drift-envelope-global}
	Let $\rho:=2d/n$ and let $B_t$ be the buffer from Lemma~\ref{lem:buffer}, so that
	$B_{t+1}-B_t=\theta-Y_{z_{t+1}}$ with $\theta=\tfrac{1}{c}-1$ and $Y_{z_{t+1}}\ge 0$ the phase extra HD.
	For any $0<\lambda<\tfrac12\ln\!\tfrac{1}{2\rho}$ define
	\[
	K(\lambda):=\frac{1-\rho}{1-\rho e^{2\lambda}}
	\qquad\text{and}\qquad
	\alpha:=e^{-\lambda\theta}\,K(\lambda).
	\]
	Then, uniformly for all $t\ge 0$,
	\[
	\EE\!\big[e^{-\lambda(B_{t+1}-B_t)}\mid \mathcal F_t\big]\ \le\ \alpha.
	\]
\end{lemma}

\begin{proof}
	By Lemma~\ref{lem:buffer}, $B_{t+1}-B_t=\theta-Y_{z_{t+1}}$, hence
	\[
	e^{-\lambda(B_{t+1}-B_t)}\ =\ e^{-\lambda\theta}\,e^{\lambda Y_{z_{t+1}}}.
	\]
	By Lemma~\ref{lem:phase-mgf} (refined phase MGF), for $0<\lambda<\tfrac12\ln\!\tfrac{1}{2\rho}$ we have the uniform envelope
	$\EE[e^{\lambda Y_{z_{t+1}}}\mid\mathcal F_t]\le K(\lambda)$ for all phases (all $z_{t+1}\le d$). Therefore
	\begin{align*}
		\EE\!\big[e^{-\lambda(B_{t+1}-B_t)}\mid \mathcal F_t\big]
		\ &=\ e^{-\lambda\theta}\,\EE\!\big[e^{\lambda Y_{z_{t+1}}}\mid\mathcal F_t\big]
		\ \\ &\le\ e^{-\lambda\theta}\,K(\lambda)
		\ =\ \alpha	
	\end{align*}
\end{proof}


\setcounter{corollary}{\numexpr\getrefnumber{cor:ville-exponential}-1\relax}

\begin{corollary}[Ville barrier bound for a process with an exponential mgf drift]
	Let $(B_t)_{t\ge 0}$ be adapted to $(\mathcal F_t)_{t\ge 0}$ with $B_0=b>0$.
	Assume there exist $\lambda>0$ and $\alpha\in(0,1)$ such that, for all $t\ge 0$,
	\[
	\EE\!\big[e^{-\lambda(B_{t+1}-B_t)}\mid \mathcal F_t\big]\ \le\ \alpha.
	\]
	Then
	\[
	\Pr\!\big(\inf_{t\ge 0} B_t \le 0\big)\ \le\ e^{-\lambda b}.
	\]
\end{corollary}

\begin{proof}
	Define $M_t:=\alpha^{-t}e^{-\lambda B_t}$. The assumption implies
	$\EE[M_{t+1}\mid\mathcal F_t]\le M_t$, so $(M_t)$ is a nonnegative supermartingale.
	Applying Ville’s inequality (Theorem~\ref{thm:ville}) with $a=1$ gives
	$\Pr(\sup_{t\ge 0} M_t\ge 1)\le \EE[M_0]=e^{-\lambda b}$.
	Since $B_t\le 0$ for some $t$ implies $M_t\ge 1$, we obtain the claim.
	
\end{proof}


%SSSSSSSSSSSSSSSSSSSSSSSSSSSSSSSSSSSSSSSSSSSSSSSSS
%%%%%%%%%%%%%%%%%%%%%%%%%%%%%%%%%%%%%%%%%%%%%%%%%%%%%%%%%%%%%%%%%%%%
%%%%%%%%%%%%%%%%%%%%%%%%%%%%%%%%%%%%%%%%%%%%%%%%%%%%%%%%%%%%%%%%%%%%
%%%%%%%%%%%%%%%%%%%%%%%%%%%%%%%%%%%%%%%%%%%%%%%%%%%%%%%%%%%%%%%%%%%%
%%%%%%%%%%%%%%%%%%%%%%%%%%%%%%%%%%%%%%%%%%%%%%%%%%%%%%%%%%%%%%%%%%%%
%%%%%%%%%%%%%%%%%%%%%%%%%%%%%%%%%%%%%%%%%%%%%%%%%%%%%%%%%%%%%%%%%%%%
%%%%%%%%%%%%%%%%%%%%%%%%%%%%%%%%%%%%%%%%%%%%%%%%%%%%%%%%%%%%%%%%%%%%

\RestateCounter{lemma}{lem:c-threshold-instance}

%\setcounter{lemma}{\numexpr\getrefnumber{}-1\relax}

\begin{lemma}[Per-instance $c$-threshold for the Ville margin]
	%\label{lem:c-threshold-instance}
	Let $\rho:=\tfrac{2d}{n}\in(0,1)$ and $\theta=\tfrac{1}{c}-1$. Define
	\[
	c_\star(n,d)\ :=\ \frac{1-\rho}{1+\rho}\ =\ \frac{1-\tfrac{2d}{n}}{1+\tfrac{2d}{n}}.
	\]
	If $c<c_\star(n,d)$, then there exists $\lambda\in\!\big(0,\tfrac12\ln\tfrac{1}{2\rho}\big)$ such that
	\[
	\alpha(\lambda)\ :=\ e^{-\lambda\theta}\,\frac{1-\rho}{\,1-\rho e^{2\lambda}\,}\ <\ 1,
	\]
	and hence the supermartingale/Ville bound applies (Cor.~\ref{cor:ville-exponential}).
\end{lemma}

\begin{proof}
	Let $g(\lambda):=\ln\alpha(\lambda)=-\lambda\theta+\ln\!\big(\tfrac{1-\rho}{1-\rho e^{2\lambda}}\big)$. 
	Then $g(0)=0$ and $g'(0)=-\theta+\tfrac{2\rho}{1-\rho}$. 
	If $\theta>\tfrac{2\rho}{1-\rho}$ (equivalently $c<\tfrac{1-\rho}{1+\rho}$), continuity gives $g(\lambda)<0$ for all sufficiently small $\lambda>0$, i.e., $\alpha(\lambda)<1$ in the admissible domain.
\end{proof}


\setcounter{lemma}{\numexpr\getrefnumber{prop:climb-time}-1\relax}

\begin{lemma}[Monotone second-slope climb time]
	On the event $\{\inf_{t\ge 0} B_t>0\}$ (no jump-back), the second-slope ascent completes in expected time
	\[
	\EE[T_{\mathrm{climb}}]\ =\ \sum_{z=1}^{d}\frac{n}{z}\ =\ \Theta(n\log d),
	\]
	and, for $\tau=\Omega(n^{1+\varepsilon})$ with fixed $\varepsilon>0$, one has $\EE[T_{\mathrm{climb}}]=o(\tau)$.
\end{lemma}

\begin{proof}
	At zero-count $z$, success probability per attempt is at least $z/n$, so the expected number of attempts for $z\!\to\! z-1$ is at most $n/z$.
	Summing over $z=1,\dots,d$ gives $\sum_{z=1}^{d}n/z=\Theta(n\log d)$.
	Since $d\le(\tfrac14-\varepsilon)n$, we have $\log d=O(\log n)$, hence $n\log d=o(n^{1+\varepsilon})$ and $\EE[T_{\mathrm{climb}}]=o(\tau)$.
\end{proof}



%\setcounter{corollary}{\numexpr\getrefnumber{cor:c-floor}-1\relax}
\begin{corollary}[Uniform floor $\ge \tfrac13$ over the admissible $d$-range]
	\label{cor:c-floor}
	If $d\le (\tfrac14-\varepsilon)n$ with fixed $\varepsilon\in(0,1/4)$, then
	\[
	c_\star(n,d)\ =\ \frac{1-\tfrac{2d}{n}}{1+\tfrac{2d}{n}}
	\ \ge\  \frac{\tfrac12+2\varepsilon}{\tfrac32-2\varepsilon}
	\ =\ \frac{1+4\varepsilon}{3-4\varepsilon}\ \ge\ \tfrac13.
	\]
	Thus the permitted constant $c$ can always be chosen at least $1/3$ (and strictly larger for any fixed $\varepsilon>0$).
\end{corollary}


\setcounter{lemma}{\numexpr\getrefnumber{lem:param-constant-c-general}-1\relax}
\begin{lemma}[Constant safety probability from an arbitrary start (Ville route)]
	
	Assume $d\le(\tfrac14-\varepsilon)n$ and fix $c\in(0,c_\star(n,d))$ with 
	$c_\star(n,d)=\frac{1-2d/n}{1+2d/n}$.
	Let $r_0,\HD(r_0)\in\mathbb{N}$ and set $b:=B_{r_0}=\frac{r_0}{c}-\HD(r_0)>0$.
	Then there exists $\lambda\in\big(0,\tfrac12\ln\frac{n}{4d}\big)$ such that
	\[
	\Pr\!\big(\inf_{t\ge 0} B_t \le 0\big)\ \le\ e^{-\lambda\,b}.
	\]
\end{lemma}
\begin{proof}
	By Lemma~\ref{lem:c-threshold-instance} (with $\rho=2d/n$), there exists
	$\lambda\in\big(0,\tfrac12\ln\tfrac{1}{2\rho}\big)=\big(0,\tfrac12\ln\tfrac{n}{4d}\big)$
	such that $e^{-\lambda\theta}K(\lambda)<1$. Applying Corollary~\ref{cor:ville-exponential}
	to $B_t$ started at $B_{r_0}=b>0$ yields
	$\Pr(\inf_{t\ge 0} B_t\le 0)\le e^{-\lambda b}$.
\end{proof}

\textcolor{blue}{Now that the supporting lemmata are established.}

\begin{proof}[Proof of Theorem~\ref{th:expoHD-cliff}]
	\emph{Up to the second slope.}
	From a uniformly random start, the expected iterations to reach the cliff-top ($n-d$ ones) are
	$\sum_{i=d}^{n}\frac{n}{i}=\Theta\!\big(n\log\tfrac{n}{d}\big)$, which at $\Theta(n)$ evaluations/iteration costs
	$\Theta\!\big(n^{2}\log\tfrac{n}{d}\big)$ evaluations per ascent to the cliff-top. At the cliff-top, the only fitter point is $1^n$, so ageing triggers unless the optimum is found within $\tau$ iterations.
	
	By Lemma~\ref{lem:can-start-omega-d3} (canonical start from an ageing cliff-top), within at most three iterations after an ageing kick the process reaches the canonical second-slope start
	\[
	r_0=2,\qquad \HD(r_0)=2
	\]
	with probability
	\[
	p_{\mathrm{can}}(n,d)\ \ge\ \frac{3}{256}\Big(\frac{d}{n}\Big)^{\!3}.
	\]
	Hence the number of ageing epochs needed to obtain this start is geometric with mean $1/p_{\mathrm{can}}$, and the expected evaluation cost up to the canonical start (including restarts) is
	\[
	\Theta\!\Big(\frac{1}{p_{\mathrm{can}}(n,d)}\cdot n^{2}\log\tfrac{n}{d}\Big)
	\ =\ O\!\Big(\frac{n^{5}\log(n/d)}{d^{3}}\Big).
	\]
	
	\smallskip
	\emph{Second–slope control (Ville route).}
	Under the reference-update rule the start on the second slope is $r_0=2$, $\HD(r_0)=2$, so $s_0=0$ and with $\theta=\tfrac1c-1$ define
	\[
	B_r\ :=\ \theta r - \sum_{j=r_0+1}^{r} Y_{z_j}\qquad(r\ge r_0),
	\]
	so $B_{r_0}=2\theta$ and $B_{t+1}-B_t=\theta-Y_{z_{t+1}}$ by Lemma~\ref{lem:buffer}.
	By the refined phase MGF (Lemma~\ref{lem:phase-mgf}), for every $0<\lambda<\tfrac12\ln\!\tfrac{1}{2\rho}$ with $\rho:=2d/n$,
	\[
	\EE\!\big[e^{\lambda Y_z}\big]\ \le\ K(\lambda)\ :=\ \frac{1-\rho}{\,1-\rho e^{2\lambda}\,}.
	\]
	Therefore, by Lemma~\ref{lem:drift-envelope-global},
	\[
	\EE\!\big[e^{-\lambda(B_{t+1}-B_t)}\mid\mathcal F_t\big]\ \le\ e^{-\lambda\theta}\,K(\lambda)\qquad\text{for all }t\ge 0.
	\]
	
	\emph{Safety bound (instantiating the per-instance lemma uniformly).}
	Since $d\le(\tfrac14-\varepsilon)n$ and the theorem assumes
	\[
	c\ <\ \frac{1+4\varepsilon}{3-4\varepsilon}\ \le\ \frac{1-2d/n}{1+2d/n}\ =:\ c_\star(n,d),
	\]
	Lemma~\ref{lem:c-threshold-instance} guarantees an admissible $\lambda\in\big(0,\tfrac12\ln\tfrac{n}{4d}\big)$ with $e^{-\lambda\theta}K(\lambda)<1$.
	With $B_{r_0}=2\theta$, Corollary~\ref{cor:ville-exponential} yields
	\[
	\Pr\!\Big(\inf_{t\ge 0} B_t\le 0\Big)\ \le\ e^{-\lambda\,2\theta}.
	\]
	Moreover, since $\tfrac12\ln\!\tfrac{n}{4d}\ge \tfrac12\ln\!\tfrac{1}{1-4\varepsilon}$ for all admissible $d$, we may fix a \emph{uniform} $\lambda_\varepsilon\in\big(0,\tfrac12\ln\tfrac{1}{1-4\varepsilon}\big)$ and conclude
	\[
	\Pr\!\Big(\inf_{t\ge 0} B_t\le 0\Big)\ \le\ e^{-\lambda_\varepsilon\,2\theta}.
	\]
	Because $\lambda_\varepsilon$ depends only on $\varepsilon$ and $\theta=\tfrac1c-1$ depends only on $c$, this right-hand side is a constant in $(0,1)$, independent of $n$ and $d$.
	
	\smallskip
	\emph{Time to finish once safe.}
	On the event $\{\inf_{t\ge 0} B_t>0\}$ (no jump-back), the second-slope ascent is monotone and completes in expected
	\[
	\sum_{z=1}^{d}\frac{n}{z}\ =\ \Theta(n\log d)
	\]
	iterations, which is $o(\tau)$ since $\tau=\Omega(n^{1+\epsilon'})$. With a pessimistic $\Theta(n)$ evaluations per iteration, one second-slope attempt costs $\Theta(n^{2}\log d)$ evaluations, and the constant no-jump-back probability implies a constant expected number of such attempts.
	
	\smallskip
	\emph{Total.}
	Summing the (restarted) approach to the canonical start and the second-slope cost gives
	\[
	\EE[\text{evaluations}]
	\ =\ O\!\Big(\frac{n^{5}\log(n/d)}{d^{3}}\Big)\ +\ O\!\big(n^{2}\log d\big),
	\]
	as claimed.
\end{proof}

